\pagebreak

\chapter{Miscellaneous tools}\label{ch:misctools}

    The 3PL source directory {\bf sources/tools} contains miscellaneous
    scripts and programs related to configuration file generation or CPU to FPGA
    interfaces. The {\bf make} files for each of these installs them in 
    {\bf \verb'~'/install/3pl}. In most cases the installation simply copies the files.
    
    The following utility programs should be executed on the host computer -
    
    \hspace*{1cm}{\bf bin2c} \\
    \hspace*{1cm}{\bf bin2py} \\
    \hspace*{1cm}{\bf info2c} \\
    \hspace*{1cm}{\bf info2py} \\
    \hspace*{1cm}{\bf edif2bit} \\
    \hspace*{1cm}{\bf run\_ISE} \\
    \hspace*{1cm}{\bf run\_Vivado} \\
    
    {\bf bin2c} and {\bf bin2py} can be run on the Zynq CPU however they require
    the .bit file as input. Similarly {\bf info2c} and {\bf info2py} can be run
    on the Zynq CPU but they require the .info file as input.

    edif2bit, run\_ISE and run\_Vivado execute Xilinx software tools and hence
    cannot execute on the Zynq CPU.

    The remaining tools are used by CPU code on the Zynq and hence should be
    copied to the Zynq CPU, probably again to {\bf \verb'~'/install/3pl} there -
    
    \hspace*{1cm}{\bf cpu\_serial.c} \\
    \hspace*{1cm}{\bf cpu\_serial.h} \\
    \hspace*{1cm}{\bf cpu\_serial.py} \\
    \hspace*{1cm}{\bf fpga\_py\_extension.c} \\
    \hspace*{1cm}{\bf fpga.c} \\
    \hspace*{1cm}{\bf fpga.h} \\
    \hspace*{1cm}{\bf fpga.py} \\
    \hspace*{1cm}{\bf mk\_setup} \\
    \hspace*{1cm}{\bf setup.py} \\
            
    \section{bin2c/}
    
        Type: Python3 script\\
        Source: {\bf sources/tools/bin2c}\\
        Target: any\\
        Execute: on host computer or on Zynq CPU
    
        Convert a binary compressed FPGA configuration bit file to a C
        include file containing an initialised byte array as C source code.
        By this means the configuration can be included in the application
        program and can be loaded by the CPU application code without the
        need to read an external file during execution.        
        
        Execute by -
        
        {\bf bin2c {\em design\_name}}
        
        which will read {\bf {\em
        design\_name}.bit.gz} and produce a C include file {\bf {\em
        design\_name}\_include.h}.
        
        This is currently only applicable to the Zynq platform. In principle it
        could be used for platforms whose configuration can be loaded via a
        serial link from CPU to FPGA, or some other channel, but a procedure to
        do that would have to be created for each specific platform.
        
        The following application code for design name {\em prog} illustrates
        the inclusion of the configuration data and the configuration
        loading procedure -
        \begin{lstlisting}[gobble=4, language=C]
            .
            .
            
            #include "prog_cinclude.h"
            .
            .

            int main (int argc, const char *argv[]) {

                fpga_open();

                fpga_configdata_load(prog_fpga_config_data, prog_fpga_config_size);

                .
                .
            }
        \end{lstlisting}
        The source for {\bf fpga\_configdata\_load} is in 
        {\bf sources/tools/fpga} and also {\bf \verb'~'/install/3pl/src/fpga.c}
        and the object file is {\bf \verb'~'/install/3pl/lib/fpga.o}.
        
        Variables {\bf prog\_fpga\_config\_data} and {\bf prog\_fpga\_config\_size}
        are defined in the included file {\bf prog\_cinclude.h}.

    \section{bin2py/}
    
        Type: Python3 script\\        
        Source: {\bf sources/tools/bin2py}\\
        Target: any\\
        Execute: on host computer or on Zynq CPU
    
        Convert a binary compressed or uncompressed FPGA configuration bit file
        to a Python3 source file containing an initialised byte array as Python3
        source code to be imported. By this means the configuration can be
        embedded in the application program and can be loaded by the CPU
        application code without the need to read an external file during
        execution.
        
        Execute by -
        
        {\bf bin2py {\em design\_name}}
        
        which will read {\bf {\em design\_name}.bit} or {\bf {\em
        design\_name}.bit.gz} and produce a Python3 file {\bf {\em
        design\_name}\_config.py}.
        
        The following application code for design name {\em prog} illustrates
        the inclusion of the configuration data and the configuration
        loading procedure -
        \begin{lstlisting}[gobble=4, language=Python]
        import sys
        import os
        sys.path.append(os.environ['HOME'] + "/install/3pl")
        from fpga import *
        import prog_fpga as fpga
        import prog_config as config

        FPGA_OPEN()
        FPGA_CONFIGDATA_LOAD(config.data, len(config.data))
        .
        .
        \end{lstlisting}
        In the above {\bf from fpga import *} defines low level functions for
        communication between the CPU and the FPGA.       
        
        {\bf import prog\_fpga as fpga} imports a definition of the application
        level FPGA interface functions and has been created by executing {\bf
        info2py}. The same effect could be obtained by importing {\bf info2py.py}.
        
        {\bf import prog\_config as config} imports a Python3 source version of the
        FPGA configuration file (preferably compressed) and is created by {\bf
        edif2bit}.


        
    \section{info2c/}

        Type: Python3 script\\
        Source: {\bf sources/tools/info2c}\\
        Target: Zynq or any with serial interface\\       
        Execute: on host computer or on Zynq CPU

        For a Zynq system, or a board with a serial CPU interface, generate
        CPU/FPGA communication functions for C from the 3PL
        information file (design.info). This generates an include file
        {\bf {\em design\_name}\_include.h}.

    \section{info2py/}
    
        Type: Python3 script        
        Source: {\bf sources/tools/info2py}\\
        Target: Zynq or any with serial interface\\       
        Execute: on host computer or on Zynq CPU
    
        Generate CPU/FPGA interface functions for
        Python3 from the 3PL information file (design.info).
        
        This can be executed as a main program (\verb'~'/install/3pl/bin/info2py) and will
        generate a Python3 source file of CPU/FPGA interface functions that can be
        imported into the CPU application program.
        
        Alternatively this file (\verb'~'/install/3pl/include/info2py.py) can be imported
        into the application program in which case it will create the interface
        functions internally and will compile them into a module.

        info2py is a copy of into2py.py included so that it can be
        invoked as a main program without the file name extension, e.g.
        {\bf info2py prog} rather than {\bf python3 infotopy.py prog}.

        This source will import {\bf serial/cpu\_serial/cpu\_serial.py}
        for FPGA platforms which have a serial interface, or
        {\bf zynq/fpga/fpga.py} in the case of Zynq.

        To execute {\bf info2py} and get an importable file -
        
        {\bf info2py {\em design\_name}}
        
        which will read {\bf {\em design\_name}.info} and create a Python3
        source file {\bf {\em design\_name}\_fpga.py} which can be imported by the
        Python3 application program. The imported file defines all the FPGA
        interface functions.
        
        If info2py.py is imported into the application program the FPGA interface
        JSON description can be read from the .info file or from a version stored
        in block RAM in the FPGA. The following code illustrates both of these
        cases with one of them commented out -
        
        \begin{lstlisting}[gobble=4, language=Python]
        import sys
        import os
        sys.path.append(os.environ['HOME'] + "/install/3pl")
        from fpga import *
        import info2py

        FPGA_OPEN()

        # From .info file
        #fpga = info2py.info_from_file("prog")
        
        # From FPGA
        fpga = info2py.info_from_fpga()
        .
        .
        \end{lstlisting}






    \section{edif2bit/}
        
        Type: shell/TCL script\\        
        Source: {\bf sources/tools/edif2bit}\\
        Target: any\\        
        Execute: on a Linux host computer
        
        Generate an FPGA configuration bit file
        from the EDIF netlist file. This invokes the Xilinx ISE or
        Vivado tools as appropriate and hence must be executed on a host
        containing these tools.
        
        Execute by -
        
        {\bf edif2bit {\em design\_name}}
        
        which will read {\bf {\em design\_name}.edif} and either
        {\bf {\em design\_name}.ncf}
        (ISE) or {\bf {\em design\_name}.xdc} (Vivado) and produce the configuration
        file {\bf {\em design\_name}.bit}.
        
        optional command line arguments include -tools ("ise" or "vivado") to
        manually select place-and-route software, -compress to run gzip on the
        bit file and -info to append the .info file and .rpt file to the
        configuration file. This last option appends meta-data to the configuration
        file to document its production.
        
        The bit file can be incorporated into the CPU application program by using
        bin2c or bin2py and then loaded into the FPGA during proggram execution as
        described in previous sections.
        
        Alternatively the FPGA configuration can be loaded from the bit file
        by the application program. The bit file may be compressed,
        i.e. "prog.bit" or "prog.bit.gz".
        
        The following example is in C -
        \begin{lstlisting}[gobble=4, language=C]
        #include <stdio.h>
        #include <stdint.h>
        #include "prog_fpga.h"

        int main (int argc, const char *argv[]) {
            int         i;

            fpga_open();
            fpga_configfile_load("prog.bit");
            .
            .
        \end{lstlisting}
        
        Or in Python3 -
        
        \begin{lstlisting}[gobble=4, language=Python]
        import sys
        import os
        sys.path.append(os.environ['HOME'] + "/install/3pl")
        from fpga import *
        import prog_fpga as fpga
        import time
        import signal
        import prog_config as config

        FPGA_OPEN()
        FPGA_CONFIGFILE_LOAD("prog.bit.gz"))
        .
        .
        \end{lstlisting}




                
    \section{run\_ISE}
    
        Type: shell script\\        
        Source: {\bf sources/tools/run\_ISE}\\
        Target: any\\
        Execute: on a Linux host computer
    
        Generate an FPGA configuration bit file
        from the EDIF netlist file using Xilinx ISE tools. This is
        a very much more primitive and direct script than edif2bit
        above.
        
        Execute by -
        
        {\bf run\_ISE {\em design}\_name}

    \section{run\_vivado}
    
        Type: shell script\\        
        Source: {\bf sources/tools/run\_Vivado}\\
        Target: any\\
        Execute: on a Linux host computer
    
        Generate an FPGA configuration bit
        file from the EDIF netlist file using Xilinx Vivado
        tools. This is a very much more primitive and direct
        script than edif2bit above.
        
        Execute by -
        
        {\bf run\_Vivado {\em design}\_name}

    \section{serial/}
    
        This directory contains interface code for FPGA platforms which have
        a serial interface between the CPU and the FPGA. An example of this
        type of platform is the Gadget Factory Papilio board.
        
        These procedures can be called directly but for convenience are usually
        called through the included/imported intervening procedures generated by
        info2c or info2py
                
        See also - \autorefph{ch:serialcpu}.
                
        
        \subsection{cpu\_serial.py}\label{ch:tools_python_serial}
    
            Type: Python3\\
            Source: {\bf sources/tools/serial/cpu\_serial.py}\\
            Target: platforms with serial interface to CPU\\
            Execute: on a Linux or OSX host computer
            
            As well as the open and close procedures listed below this file
            defines necessary low-level CPU/FPGA interface communications functions.
            These functions are called from the interface code in the Python3 file
            generated by {\bf info2py} and imported into the application code.
            
            \underline{{\bf fpga\_open(port\_name)}} \\
            Open the serial interface to the FPGA. {\bf port\_name} is a string
            specifying the serial device to the FPGA, e.g.
            "/dev/tty.usbmodem14101".
            
            \underline{{\bf fpga\_close()}} \\
            Close the serial interface to the FPGA.
            
            \underline{{\bf DMA\_buffer(b)}} \\
            Set a pointer to the DMA buffer, eg.
            
            \begin{lstlisting}[gobble=4, language=Python]
            dmabuffer = bytearray(4096)
            DMA_buffer(dmabuffer)
            \end{lstlisting}

            
                
        \subsection{cpu\_serial.c/, cpu\_serial.h}\label{ch:tools_c_serial}
    
            Type: C\\
            Source: {\bf sources/tools/serial/cpu\_serial.c}, {\bf sources/tools/serial/cpu\_serial.h} \\
            Target: platforms with serial interface to CPU\\
            Execute: on a Linux or OSX host computer
        
            
            As well as the procedures listed below this file defines necessary
            low-level CPU/FPGA interface communications functions. These
            functions are called from the procedures below and from interface
            code in the C include file generated by {\bf info2c}.
            
            \underline{{\bf  fpga\_open(char* port\_name))}} \\
            Open the serial interface to the FPGA. {\bf port\_name} is a string
            specifying the serial device to the FPGA, e.g.
            "/dev/tty.usbmodem14101".
            
            \underline{{\bf fpga\_close()}} \\
            Close the serial interface to the FPGA.
            
            \underline{{\bf set\_dma\_address(char* addr))}} \\
            Set a DMA base address for the pseudo-DMA.

        
            

    \section{zynq/}
    
        This directory contains CPU/FPGA interface functions for the Zynq FPGA
        platform.
        
        These procedures can be called directly but for convenience are usually
        called through the included/imported intervening procedures generated by
        info2c or info2py
        
        See also - \autorefph{ch:zynqcpu}.
    
        For this interface to operate the Zynq Linux device tree must provide
        entries for axi\_mgp0, axi\_mgp1 and axi\_hpbuf and must also provide
        interrupt events 0 to 15. An example is shown in
        \autorefph{ch:zynq_device_tree}.
            
        
        \subsection{fpga.c, fpga.h}\label{ch:tools_c_zynq}

            Type: C source\\        
            Source: {\bf sources/tools/zynq/fpga/fpga.c}, {\bf sources/tools/fpga/fpga.h}\\
            Target: Zynq\\
            Execute: on Zynq CPU

            \underline{{\bf fpga\_open ()}} \\

                This maps register space (AXI mgp0), memory space (AXI mgp1), DMA space (AXI HP0, HP1, HP2 and
                HP3) and events 0 to 15.

            \underline{{\bf fpga\_close ()}} \\

                This unmaps the previously mapped spaces.

            \underline{{\bf fpga\_event\_wait(uint32\_t event\_index)}} \\

                Wait for an event. This enables the event specifies by the event index {\bf ev}
                and then does a blocking read on /dev/uiox where x is the event unit number. When
                the read occurs the event is now disabled.

            \underline{{\bf fpga\_event\_control(uint32\_t event\_index, uint32\_t c)}} \\

                Enable or disable the event specified by event index {\bf ev}. This is done by
                writing 1 or 0 respectively to /dev/uiox where x is the event unit number.

            \underline{{\bf uio\_event\_fd(uint32\_t event\_index)}} \\

                Return the file descriptor for /dev/uiox where x is the event unit number.

            \underline{{\bf uio\_event\_count(uint32\_t event\_index)}} \\

                Return the number of events since the previous wait. If greater than 1 then
                a number of events have occured before the wait has executed iyts blocking read.

            \underline{{\bf uio\_dma\_phys\_address()}} \\

                Return the physical address of the mapped DMA space.

            \underline{{\bf fpga\_configfile\_load(char *filename)}} \\

                Load an FPGA configuration from the specified bit file.

            \underline{{\bf fpga\_configdata\_load(uint8\_t *data, size\_t len)}} \\

                Load an FPGA configuration from an array of char.\\
                {\bf data} is the array of char.\\
                {\bf len} is the dimension of the array.

            The following macros are defined in fpga.h and are intended to
            be accessed via functions and procedures defined in the include
            file generated by info2c rather than being called directly from
            application code.

            \underline{{\bf FPGA\_REG\_READ(addr, type, dww)}} \\

                Read an FPGA value or queue mode variable.\\
                {\bf addr} is the address of the interface.\\
                {\bf type} is the type of the returned value.\\
                {\bf dww} is the number of 32-bit words to be transfered.

            \underline{{\bf FPGA\_REG\_WRITE(addr, val, dww)}} \\

                Write to an FPGA static or queue mode variable.\\
                {\bf addr} is the address of the interface.\\
                {\bf val} is the value to be written.\\
                {\bf dww} is the number of 32-bit words to be transfered.

            \underline{{\bf FPGA\_MEM\_READ(addr, loc, type, dww)}} \\

                Read an FPGA cmemory or rmemory mode variable.\\
                {\bf addr} is the address of the interface.\\
                {\bf loc} is the address within the memory.\\
                {\bf type} is the type of the returned value.\\
                {\bf dww} is the number of 32-bit words transfered.

            \underline{{\bf FPGA\_MEM\_WRITE(addr, loc, val, dww)}} \\

                Write to an FPGA cmemory or rmemory mode variable.\\
                {\bf addr} is the address of the interface.\\
                {\bf loc} is the address within the memory.\\
                {\bf val} is the value to be written.\\
                {\bf dww} is the number of 32-bit words transfered.


        \subsection{fpga.py}\label{ch:tools_python_zynq}

            Type: Python3 source\\        
            Source: {\bf sources/tools/zynq/fpga/fpga.py}\\
            Target: Zynq\\
            Execute: on Zynq CPU

            Python3 low-level CPU/FPGA interface communications. functions. This
            imports low-level communication functions from {\bf fpga\_read\_write.c}
            which has been added to the Python3 library by {\bf
            sources/tools/zynq/lib\_py\_ext/setup.py} below as Python3 module {\bf
            fpga.py}.

            \underline{{\bf fpga\_open ()}} \\

                This calls fpga\_c\_open() in fpga\_read\_write.c (below) to map register space
                (AXI mgp0), memory space (AXI mgp1), DMA space (AXI HP0, HP1, HP2 and HP3) and
                events 0 to 15.

                It then saves the current value of each event device /sys/class/uio/uiox/event where x is 0 to
                15. This value increments on each event occurence since oper-up and the initial value prior to
                execution is needed to determine how many events have occured between interrupts (blocking reads
                of the event device /dev/uiox).

            \underline{{\bf fpga\_close ()}} \\

                This calls fpga\_c\_close() in fpga\_read\_write.c (below) to unmap
                the previously mapped spaces.

            \underline{{\bf event\_wait (ev)}} \\

                Wait for an event. This enables the event specifies by the event index {\bf ev}
                and then does a blocking read on /dev/uiox where x is the event index. When
                the read occurs the event is now disabled.

            \underline{{\bf event\_enable (ev)}} \\

                Enable the event specified by event index {\bf ev}. This is done by
                writing 1 to /dev/uiox where x is the event index.

            \underline{{\bf event\_disable (ev)}} \\

                Disable the event specified by event index {\bf ev}. This is done by
                writing 0 to /dev/uiox where x is the event index.

            \underline{{\bf dma\_phys\_address ()}} \\

                This function returns the physical address of the DMA mapped space.

            \underline{{\bf dma\_get\_word (addr)}} \\

                This function gets one 32-bit word from the DMA mapped space. {\bf addr}
                is an address offset into the DMA address space.

            \underline{{\bf dma\_put\_word (addr, data)}} \\

                This function puts one 32-bit word into the DMA mapped space. {\bf addr}
                is an address offset into the DMA address space and {\bf data} is the data word.


            The following are called via a module generated by info2py and are
            not intended to be called directly from application code.

            \underline{{\bf comms\_write (mem, address, memaddress, dba)}} \\

                Write to the FPGA.\\
                {\bf mem} is 0 for a register or queue, 1 for cmemory or rmemory.\\
                {\bf address} is the address of the interface.\\
                {\bf memaddress} is the address within a memory if {\bf mem} is 1.\\
                {\bf dba} is a byte array containing the data. The size of the byte array
                indicates the number of bytes to write.

            \underline{{\bf comms\_read (mem, address, memaddress, bytes)}} \\

                Read from the FPGA.\\
                {\bf mem} is 0 for a value or queue, 1 for cmemory or rmemory.\\
                {\bf address} is the address of the interface.\\
                {\bf memaddress} is the address within a memory if {\bf mem} is 1.\\
                {\bf bytes} is the number of bytes to read.\\
                The data is returned as a byte array.

        \subsection{fpga\_py\_extension.c}

            Type: C source\\        
            Source: {\bf sources/tools/zynq/fpga\_py\_extension.c}\\
            Target: Zynq\\
            Execute: on Zynq CPU

            Python3 C extension code. This provides links from the Python3 interface functions
            into functions in {\bf fpga.c}. It is added to the Python3 library as module
            {\bf fpga\_read\_write} by executing {\bf setup.py} (below).


            \underline{{\bf py\_fpga\_open()}} \\

                Called as fpga\_open() in python.\\
                Calls fpga\_open() in fpga.c.
                This maps register space (AXI mgp0), memory space (AXI mgp1), DMA space (AXI HP0, HP1, HP2 and
                HP3) and events 0 to 15.

            \underline{{\bf py\_fpga\_close()}} \\

                Called as fpga\_close() in python.\\
                Calls fpga\_close() in fpga.c.
                This unmaps the previously mapped spaces.

            \underline{{\bf py\_fpga\_configfile\_load(fname)}} \\

                Called as fpga\_configfile\_load() in python.\\
                Calls fpga\_configfile\_load() in fpga.c.\\
                {\bf fname} is the file name of the configuration bit file. This will have a .bit
                file name extension or .bit.gz if compressed,

            \underline{{\bf py\_fpga\_configdata\_load(data, len)}} \\

                Called as fpga\_configdata\_load(data, len) in python.\\
                {\bf data} is a byte array containing the configuration bit pattern.\\
                {\bf len} is the dimension of the array.\\
                The configuration byte array is contained in an imported module created by {\bf bin2py}
                (which itself is a Python3 script).

            \underline{{\bf py\_fpga\_burst\_read(is\_mem, address, n\_words)}} \\

                Called as fpga\_burst\_read(mem, adress, words) in python.\\
                Read an FPGA value, queue, cmemory or rmemory mode variable.\\
                {\bf is\_mem} is 1 for cmemory or rmemory mode and 1 otherwise.\\
                {\bf address} is the address of the interface.\\
                {\bf n\_words} is the number of 32-bit words to be transfered.

            \underline{{\bf py\_fpga\_burst\_write(is\_mem, address, n\_words, data)}} \\

                Called as fpga\_burst\_write(mem, address, words, data) in python.\\
                Write to an FPGA static, queue, cmemory or rmemory mode variable.\\
                {\bf is\_mem} is 1 for cmemory or rmemory mode and 1 otherwise.\\
                {\bf address} is the address of the interface.\\
                {\bf n\_words} is the number of 32-bit words to be transfered.\\
                {\bf data} is the data, a byte array.

            \underline{{\bf py\_uio\_event\_wait(event\_no)}} \\

                Called by uio\_event\_wait(ev) in python.\\
                Calls fpga\_event\_wait() in fpga.c.\\
                {\bf event\_no} is the event index, 0 to 15.\\
                Wait for an event. This enables the event specifies by the event index {\bf ev}
                and then does a blocking read on /dev/uiox where x is the event unit number. When
                the read occurs the event is now disabled.

            \underline{{\bf py\_uio\_event\_fd(event\_no)}} \\

                Called by uio\_event\_fd(ev) in python.\\
                {\bf event\_no} is the event index, 0 to 15.\\
                Return the file descriptor for /dev/uiox where x is the event unit number.

            \underline{{\bf py\_uio\_event\_unit(event\_no)}} \\

                Called by uio\_event\_unit(ev) in python.\\
                {\bf event\_no} is the event index, 0 to 15.\\
                Return the event number x for /dev/uiox where x is the event unit number.

            \underline{{\bf py\_uio\_event\_count(event\_no)}} \\

                Called by uio\_event\_count(ev) in python.\\
                Return the number of events since the previous wait. If greater than 1 then
                a number of events have occured before the wait has executed iyts blocking read.

            \underline{{\bf py\_uio\_event\_control(event\_no, c)}} \\

                Called by uio\_event\_control(ev, c) in python.\\
                Calls fpga\_event\_control() in fpga.c.\\
                Enable or disable the event specified by event index {\bf ev}. This is done by
                writing 1 or 0 respectively to /dev/uiox where x is the event unit number.

            \underline{{\bf py\_uio\_dma\_get\_word(address)}} \\

                Called by uio\_dma\_get\_word(addr) in python.\\
                Returns a 32-bit word from the DMA space with offset {\bf address}.

            \underline{{\bf py\_uio\_dma\_put\_word(address, data)}} \\

                Called by uio\_dma\_put\_word(addr, data) in python.\\
                Copies a 32-bit word {\bf data} from the DMA space with offset {\bf address}.

            \underline{{\bf py\_uio\_dma\_phys\_address()}} \\

                Called by uio\_dma\_phys\_address() in python.\\
                Returns the physical address of the DMA space.



        \subsection{setup.py}

            Type: Python3 source\\        
            Source: {\bf sources/tools/zynq/setup.py}\\
            Target: Zynq\\
            Execute: on Zynq CPU

            A Python3 source file to add {\bf fpga\_read\_write.c} above as a C
            extension to the Python3 library. {\bf mk\_setup} is a shell
            script to execute {\bf setup.py} which will add module {\bf
            fpga\_read\_write} to the Python3 library.

        \subsection{mk\_setup}

            Type: shell command file\\        
            Source: {\bf sources/tools/zynq/mk\_setup}\\
            Target: Zynq\\
            Execute: on Zynq CPU

            This command file adds the C extension described above to the Python3 library.
            It also compiles fpga.c to fpga.o.
            Execute by typing {\bf mk\_setup} to the shell.

